\section{Words with prescribed relations and recognizable boundaries}
\label{sec:words}

We now construct finite families of words that will define the symbolic
system.  At each level, short lists of constituent labels have large
entropy, although specified pairs of labels satisfy a bilinear relation.
We also require that a word can be recognized, together with its alignment,
from a sufficiently accurate copy.  The construction in this section is
entirely finite at each step; no invariant probability is chosen yet.

Nested-word constructions and probabilistic control of frequencies and
shifted comparisons have precedents in
\cite[\S4.1 and Proposition~43]{ForemanRudolphWeiss2011}.
Related word and phase frameworks appear in
\cite[\S\S4.2.3--4.2.4]{ForemanWeiss2022}. Our graph constraints and the
per-word condition below are proved separately; in particular, we do
not require every lower-level word to occur in every parent or import
unique ergodicity from those constructions.

All positions in a finite word are numbered from zero.  For a word $u$,
write $u[t,t+q)$ for its consecutive subword of length $q$ starting at $t$.
For two words of the same length, $d_{\mathrm H}(u,v)$ denotes the number
of coordinates at which they differ.  Thus this distance counts symbols
or constituent labels according to the alphabet being used.  For
probabilities $p,q$ on a finite set, we use
\[
  \TV(p,q)=\frac12\sum_a |p(a)-q(a)|.
\]
Write $\mathcal L(Y)$ for the distribution of a random variable $Y$.

\subsection{Parameters and the four word properties}

Each new word will concatenate words from the preceding family.
We control two sampling experiments: choosing a word uniformly from
its family, and fixing a parent at the next level and choosing one of
its constituents uniformly. We compare the laws of windows at the
same internal offset in these two chosen words. Requiring this
approximation for every parent will allow arbitrary parent weights
in the eventual invariant measure.

The construction first samples constituent lists satisfying the
relations inside independent groups. It then keeps lists satisfying
the per-parent approximation, and finally selects a family with the
required window distributions and distances. The parameters below
provide enough groups and candidate words for these selections.

Fix
\begin{equation}\label{eq:word-parameters}
  \varepsilon_i=2^{-i-2},\qquad
  L_i=2^{2s_i}-1,\qquad \ell_i=\log L_i,\qquad
  K_i=\lfloor\sqrt{\ell_i}\rfloor,
\end{equation}
and put
\begin{equation}\label{eq:word-lengths}
  N_i=2^{K_i},\qquad H_i=2^{\lfloor K_i/2\rfloor},\qquad
  D_i^*=2^{\lfloor K_i/8\rfloor},\qquad
  m_1=1,\qquad m_{i+1}=m_iN_i.
\end{equation}
Starting from the integer $s_1$, define $2s_{i+1}$ to be the greatest even
integer not exceeding
\begin{equation}\label{eq:word-recursion}
  X_i:=\frac{\varepsilon_i}{100}N_i\ell_i.
\end{equation}
The proof below shows that this recursion is well defined at every level
if $K_1\ge256$.  There is no need to optimize this initial size.

Let $\mathcal A_i=\F_2^{2s_i}\setminus\{0\}$.  On the full vector space
use the nondegenerate alternating form
\[
  [u,v]_i=\sum_{a=1}^{s_i}
     \bigl(u_a v_{s_i+a}+u_{s_i+a}v_a\bigr)\in\F_2.
\]
The fixed symbol alphabet is $\mathcal A=\mathcal A_1$.
Each family $\mathcal W_i\subset\mathcal A^{m_i}$ will have $L_i$ members.
Once it has been constructed, choose a bijection with $\mathcal A_i$;
this equips its word labels with the form $[\cdot,\cdot]_i$.
In particular, a word label and its vector label are identified only
through this fixed bijection.

Whenever $w$ is a concatenation of $N_i$ members of $\mathcal W_i$, let
\[
  c_i(w)=(u_0,\ldots,u_{N_i-1})\in\mathcal W_i^{N_i}
\]
be its list of constituents.  The list is unambiguous because all members
of $\mathcal W_i$ have the same length and are distinct.
Here $m_i$ is the symbol length, $L_i$ counts distinct level-$i$
words, and $N_i$ counts constituent occurrences in a level-$(i+1)$
word, with repetitions allowed. In the properties below, $H_i$ is
the group length and $D_i^*$ bounds the window lengths and prescribed
power-of-two lags; these lengths and lags are measured in level-$i$
constituent positions.
We fix the approximation tolerance
\begin{equation}\label{eq:word-tv-tolerance}
  \vartheta=10^{-4}.
\end{equation}

\begin{proposition}[Finite word families]\label{prop:words}
Choose $s_1$ so that $K_1\ge256$, and define the parameters by
\eqref{eq:word-parameters}--\eqref{eq:word-recursion}.
There are families $\mathcal W_i\subset\mathcal A^{m_i}$ of cardinality
$L_i$, with $\mathcal W_1=\mathcal A$, such that every
$w\in\mathcal W_{i+1}$ is a concatenation of $N_i$ members of
$\mathcal W_i$ and the following properties hold.
\begin{enumerate}[label=\textup{(W\arabic*)},ref=W\arabic*,leftmargin=*]
\item\label{word:graph}\textbf{Relations within groups.}
Write $c_i(w)=(u_0,\ldots,u_{N_i-1})$.
For every $0\le b<N_i/H_i$, every
$0\le k\le\lfloor K_i/8\rfloor$, and every
$0\le r<r+2^k<H_i$,
\[
  [u_{bH_i+r},u_{bH_i+r+2^k}]_i=0.
\]
\item\label{word:distance}\textbf{Distance between words and shifted words.}
If $v,w\in\mathcal W_{i+1}$ are distinct, then
\[
  d_{\mathrm H}\bigl(c_i(v),c_i(w)\bigr)
     \ge(1-\varepsilon_i)N_i.
\]
For all $v,w,z\in\mathcal W_{i+1}$, with repetitions allowed, and all
$1\le t<N_i$,
\[
  d_{\mathrm H}\bigl((c_i(v)c_i(w))[t,t+N_i),c_i(z)\bigr)
     \ge(1-\varepsilon_i)N_i.
\]
\item\label{word:window}\textbf{Distributions of short windows.}
Let $U$ be uniform on $\mathcal W_{i+1}$.
For each $0\le t\le N_i-D_i^*$ there is a probability $Q_{i,t}$ on
$\mathcal W_i^{D_i^*}$ such that
\[
  \TV\bigl(\mathcal L(c_i(U)[t,t+D_i^*)),Q_{i,t}\bigr)
     \le\vartheta.
\]
Moreover, for $1\le d'\le D_i^*$, any distinct positions
$0\le r_1<\cdots<r_{d'}<D_i^*$, and any prescribed labels
$v_1,\ldots,v_{d'}\in\mathcal W_i$,
\begin{equation}\label{eq:word-sublist}
  Q_{i,t}\{u:u_{r_j}=v_j\text{ for }1\le j\le d'\}
       \le 2\left(\frac{2^{2K_i}}{L_i}\right)^{d'}.
\end{equation}
\item\label{word:frequency}\textbf{Approximation in each individual word.}
Suppose $i\ge2$.  For every fixed $w\in\mathcal W_{i+1}$, write
$c_i(w)=(u_0,\ldots,u_{N_i-1})$ and let $J$ be uniform on
$\{0,\ldots,N_i-1\}$.  If $U$ is uniform on $\mathcal W_i$, then for
each $0\le t\le N_{i-1}-D_{i-1}^*$,
\[
 \TV\left(
   \mathcal L\bigl(c_{i-1}(u_J)[t,t+D_{i-1}^*)\bigr),
   \mathcal L\bigl(c_{i-1}(U)[t,t+D_{i-1}^*)\bigr)
 \right)\le\vartheta.
\]
The first law averages the constituents of the single fixed word $w$;
it does not average over $\mathcal W_{i+1}$.
\end{enumerate}
\end{proposition}

Property~(\ref{word:frequency}) compares the two sampling experiments
described above. The following elementary estimate controls empirical
laws both within each parent and across the selected family.

\begin{lemma}[Averages of probability vectors]\label{lem:word-tv}
Let $Y_1,\ldots,Y_q$ be independent random probability vectors on a set
of cardinality $A$, all with mean $p$.  Then
\[
  \E\,\TV\left(\frac1q\sum_{b=1}^qY_b,p\right)
     \le\frac12\sqrt{\frac Aq}.
\]
Coordinates within an individual vector need not be independent.
\end{lemma}
\begin{proof}
For $\widehat p=q^{-1}\sum_bY_b$, independence between vectors gives
\[
 \sum_a\operatorname{Var}(\widehat p(a))
   =\frac1{q^2}\sum_{b,a}\operatorname{Var}(Y_b(a))
   \le\frac1{q^2}\sum_b\E\sum_aY_b(a)^2
   \le\frac1q.
\]
Apply the inequality $\E|Z|\le(\E|Z|^2)^{1/2}$ coordinatewise and then
Cauchy--Schwarz to the sum over the $A$ coordinates.
\end{proof}

\begin{proof}[Proof of \Cref{prop:words}]
\emph{Growth of the fixed parameters.}
Whenever $X_i\ge6$, rounding in \eqref{eq:word-recursion} gives
\begin{equation}\label{eq:word-rounding}
  X_i-3<\ell_{i+1}\le X_i,\qquad
  \ell_{i+1}\ge\frac{X_i}{2}
       =\frac{\varepsilon_i}{200}2^{K_i}\ell_i.
\end{equation}
Indeed, $X_i-2<2s_{i+1}\le X_i$, and
$2s_{i+1}-1<\log(2^{2s_{i+1}}-1)<2s_{i+1}$ when $2s_{i+1}\ge2$.
We claim inductively that
\begin{equation}\label{eq:word-growth}
 K_i\ge256,\qquad
 \varepsilon_iK_i\ge32\,2^{i-1},\qquad
 K_{i+1}\ge2^{K_i/3}\ge K_i^2\ge4K_i.
\end{equation}
The first two inequalities hold at $i=1$.  If they hold at level $i$,
then $\ell_i\ge K_i^2$ implies
\[
  X_i\ge\frac8{25}K_i2^{K_i}\ge6,
  \qquad
  \ell_{i+1}\ge\frac4{25}K_i2^{K_i}\ge2^{K_i}.
\]
Consequently
$K_{i+1}\ge2^{K_i/2}-1\ge2^{K_i/3}\ge K_i^2\ge4K_i$.
These elementary inequalities hold for every $K_i\ge256$;
for example, $K/3-2\log K$ is increasing there and is positive at $256$.
Since $\varepsilon_{i+1}=\varepsilon_i/2$, the second bound in
\eqref{eq:word-growth} propagates as well.  In particular, all parameters
are positive and defined, $H_i$ divides $N_i$, and $D_i^*\le H_i$.

\smallskip\noindent
\emph{Sampling the relations.}
Suppose that $\mathcal W_i$ and its vector labeling have been fixed.
Generate a list $(U_0,\ldots,U_{N_i-1})$ in independent consecutive
groups of $H_i$ positions.  Within each group, proceed from left to
right.  Choose each new label uniformly among the nonzero vectors
orthogonal to all earlier labels at the lags
$2^k$, $0\le k\le\lfloor K_i/8\rfloor$, that stay inside the group.  There are at most $\lfloor K_i/8\rfloor+1\le K_i$
linear equations.  Their common kernel has codimension at most $K_i$,
whether or not the equations are independent.  As $2s_i\ge\ell_i\ge K_i^2$,
the number of nonzero choices is at least
\begin{equation}\label{eq:word-choices}
  2^{2s_i-K_i}-1
    \ge2^{2s_i-K_i-1}
    \ge\frac{L_i}{2^{2K_i}}.
\end{equation}
Denote this unconditioned law on $\mathcal W_i^{N_i}$ by $P_i^0$.
Every sampled list satisfies the relations in Property~(\ref{word:graph}).

For later use, set $\zeta_i=2^{2K_i}/L_i$ within this proof.
The conditional probability of any specified label at any position,
given all preceding positions, is at most $\zeta_i$.
If $r_1<\cdots<r_{d'}$ are specified positions, conditioning first at
$r_{d'}$ and then successively at the earlier specified positions gives
\begin{equation}\label{eq:word-raw-sublist}
  P_i^0\{U_{r_j}=v_j\text{ for all }j\}\le\zeta_i^{d'}.
\end{equation}
All unspecified labels are integrated out.  This argument does not
require independence within a group.

Each individual $U_r$ nevertheless has a uniform marginal on
$\mathcal W_i$.  To see this, the group of linear maps preserving the
alternating form acts transitively on its nonzero vectors.  Given such
a vector $v$, nondegeneracy provides $w$ with $[v,w]_i=1$; the plane
spanned by $v,w$ is nondegenerate, and its orthogonal complement carries
a nondegenerate alternating form of dimension two less.  Induction
extends $v,w$ to a symplectic basis, proving transitivity by mapping one
such basis to another.  The sequential rule is equivariant under these
maps, including its uniform choices, so each marginal is invariant
under a transitive action and hence uniform.

\smallskip\noindent
\emph{Imposing the approximation in each word.}
For $i\ge2$ and an offset $0\le t\le N_{i-1}-D_{i-1}^*$, let
\[
 f_t(u)=c_{i-1}(u)[t,t+D_{i-1}^*),\qquad u\in\mathcal W_i,
\]
and let $p_t$ be its law under uniform choice of $u$.
For each of the $q_i=N_i/H_i$ independently sampled groups, form the
empirical probability vector of the $H_i$ values $f_t(U_r)$ in that
group.  Its mean is $p_t$ by the uniform one-label marginals just proved.
The average of these independent vectors is exactly the law of
$f_t(U_J)$ when $J$ is uniform over the entire sampled list.
The alphabet of $f_t$ has cardinality at most
$A_i=L_{i-1}^{D_{i-1}^*}$.
Thus \Cref{lem:word-tv}, Markov's inequality, and a union over at most
$N_{i-1}$ offsets bound the probability of violating any requirement in
Property~(\ref{word:frequency}) by
\begin{equation}\label{eq:word-filter-budget}
  F_i:=\vartheta^{-1}N_{i-1}
       \sqrt{\frac{L_{i-1}^{D_{i-1}^*}}{N_i/H_i}}.
\end{equation}
Below we verify $F_i<1/2$ at every level $i\ge2$.

Let $\mathcal E_i$ be the event that all these requirements hold, and
let $P_i=P_i^0(\,\cdot\mid\mathcal E_i)$.
For $i=1$, take $\mathcal E_1$ to be the whole sample space and
$P_1=P_1^0$.  In all cases $P_i^0(\mathcal E_i)\ge1/2$, so
\eqref{eq:word-raw-sublist} yields
\begin{equation}\label{eq:word-filtered-sublist}
 P_i\{U_{r_j}=v_j\text{ for all }j\}\le2\zeta_i^{d'}.
\end{equation}
The factor two is paid once for the whole conditioning event.
No uniform marginal or independence between groups is asserted for
$P_i$.  The relations hold on its support, as does
Property~(\ref{word:frequency}) for each individual list.

\smallskip\noindent
\emph{Selecting the finite family.}
Put $M_i=L_{i+1}$ and sample $M_i$ independent complete lists with law
$P_i$.  For each offset $0\le t\le N_i-D_i^*$, let $Q_{i,t}$ be the
law under $P_i$ of the length-$D_i^*$ window at $t$.
It satisfies \eqref{eq:word-sublist} by
\eqref{eq:word-filtered-sublist}.
The empirical window law of the $M_i$ samples is an average of
independent point-mass probability vectors on an alphabet of size
$L_i^{D_i^*}$.  A second application of \Cref{lem:word-tv} and Markov's
inequality, now with a union over at most $N_i$ offsets, bounds the
probability of violating any required window approximation by
\begin{equation}\label{eq:word-window-budget}
  E_i:=\vartheta^{-1}N_i
        \sqrt{\frac{L_i^{D_i^*}}{L_{i+1}}}.
\end{equation}
We will have $E_i<1/4$ at every level.

It remains to impose the distances, including comparisons in which
some of the sampled-word indices coincide.  First use the reference
law in which the $N_i$ labels in each distinct indexed list are
independent and uniform.  Consider one pairwise comparison between
distinct indices, or one comparison of a target list with a shifted
window from two source lists.  In the latter case write $t$ for its
offset, with $1\le t<N_i$; any of the three indices may coincide.

For a specified set of matching positions, form a multigraph whose
vertices are variables (sample index, position) and whose edges are
the asserted equalities.  Every target position occurs once.  The
source positions also occur once: they use $t,\ldots,N_i-1$ in the
first source and $0,\ldots,t-1$ in the second, disjoint ranges even if
the source indices agree.  The degree is therefore at most two,
counting repeated edges.  There is no loop.  In a pairwise comparison
the indices are distinct, and in a shifted comparison the source
position differs from the target position by $t$ modulo $N_i$.

A connected component with $v$ vertices requires $v-1$ independent
equalities and has probability $L_i^{-(v-1)}$ under the reference law.
A path has one fewer edge than vertices; a cycle has equally many,
with at least two edges.  Thus every component imposes at least half
as many independent equalities as edges, including a pair of parallel
edges.  If $e$ positions are specified to match, their probability is
at most $L_i^{-e/2}$.
Failure of a distance requirement entails at least
$\lceil\varepsilon_iN_i\rceil$ matching positions.  Summing over their
possible sets, of which there are at most $2^{N_i}$, gives the bound
\begin{equation}\label{eq:word-reference-distance}
  2^{N_i}L_i^{-\varepsilon_iN_i/2}.
\end{equation}

By \eqref{eq:word-choices}, a complete list under $P_i^0$ has likelihood
ratio at most $2^{2K_iN_i}$ relative to the reference law, and under
$P_i$ at most $2\,2^{2K_iN_i}$.
Each comparison uses at most three distinct sampled lists, so its
probability under the actual sampling law is at most
$(2\,2^{2K_iN_i})^3$ times \eqref{eq:word-reference-distance}.
Repeated indices do not introduce an additional independent list.
There are fewer than $4N_iM_i^3$ pairwise and shifted comparisons in
total.  Their union failure probability is consequently at most
\begin{equation}\label{eq:word-distance-budget}
  T_i:=32N_iL_{i+1}^3\,
       2^{(6K_i+1)N_i}L_i^{-\varepsilon_iN_i/2}.
\end{equation}
We will verify $T_i<1/4$.

These estimates address separate requirements: the first sampling law
enforces the relations, conditioning enforces the approximation in
each individual word, and the final independent samples provide the
window laws and distances.  We now check that the fixed initial size
makes all three probability bounds hold at every step.

\smallskip\noindent
\emph{Uniform bounds from one initial choice.}
For $K\ge256$, set
\[
  B_0(K)=14+K+\tfrac12\,2^{K/8}(K+1)^2.
\]
The elementary inequalities
$14+K\le2^{K/8}$ and $(K+1)^2\le2^{K/8}$ hold at $256$ and persist
thereafter, as the ratios of their left sides to their right sides
decrease.  Hence
\begin{equation}\label{eq:word-common-bound}
 B_0(K)\le2^{K/4}\le\tfrac18\,2^{K/3}.
\end{equation}
We also have $\log(\vartheta^{-1})<14$ and
$K_i^2\le\ell_i<(K_i+1)^2$.
For the first probability bound, put $K=K_{i-1}$ and use
$\log(N_i/H_i)\ge K_i/2$ to obtain
\[
 \log F_i
   \le B_0(K)-\frac{K_i}{4}
   \le-\frac18\,2^{K/3}.
\]
For the second, put $K=K_i$.  Since $\ell_{i+1}\ge2^K$,
\[
 \log E_i
   \le B_0(K)-\frac{\ell_{i+1}}2
   \le2^{K/4}-\frac{2^K}{2}
   \le-\frac{2^K}{4}.
\]
Finally \eqref{eq:word-rounding} and
$\varepsilon_i\ell_i\ge\varepsilon_iK_i^2\ge32K_i$ give
\begin{align*}
 \log T_i
 &\le5+K_i+N_i
       \left(6K_i+1-\frac{47}{100}\varepsilon_i\ell_i\right)\\
 &\le5+K_i-9K_i2^{K_i}
  \le-K_i2^{K_i}.
\end{align*}
It follows that
\begin{equation}\label{eq:word-uniform-budgets}
 F_i\le2^{-2^{K_{i-1}/3}/8}<\tfrac12\quad(i\ge2),\qquad
 E_i\le2^{-2^{K_i}/4}<\tfrac14,\qquad
 T_i\le2^{-K_i2^{K_i}}<\tfrac14.
\end{equation}
All right sides are decreasing functions of their $K$ argument.
They apply simultaneously at every level by \eqref{eq:word-growth};
no later adjustment of $s_1$ is made.

The conditioning event thus has positive probability.  Under the
conditioned law the two subsequent failures have total probability
less than $1/2$, so a tuple satisfying every requirement exists.
Its pairwise distance requirement is imposed on distinct sample
indices and has a positive lower bound.  Equal sampled words would
have distance zero, so the tuple has exactly $L_{i+1}$ distinct
members.  Flatten their constituent lists to form $\mathcal W_{i+1}$.
The empirical distribution is now exactly the uniform distribution
on this set, without deduplication or a change of weights.
It therefore has all four asserted properties.  Choose its bijection
with $\mathcal A_{i+1}$ and continue the induction.
\end{proof}

\subsection{Recognition at every symbol offset}

The distances in \Cref{prop:words} concern constituent labels.
The next consequence converts them into distances in the fixed
alphabet $\mathcal A$, including shifts that do not respect any
lower-level boundary. In coding terminology the two conclusions below
are a positive minimum distance and a positive \emph{comma-free index}:
the latter measures distance from every nontrivial splice of two legal
words. This is the synchronization notion described in
\cite[\S1]{Tonchev2003}; the recursive lower bound here is proved directly.

\begin{lemma}[Recognition of words and alignment]\label{lem:recognition}
For the families in \Cref{prop:words}, put
\[
 d_i=\prod_{j=1}^{i-1}(1-\varepsilon_j),\qquad
 d_*=\prod_{j=1}^{\infty}(1-\varepsilon_j)>0.
\]
Distinct $v,w\in\mathcal W_i$ satisfy
\[
 d_{\mathrm H}(v,w)\ge d_i m_i\ge d_*m_i.
\]
For all $v,w,z\in\mathcal W_i$ and every $1\le t<m_i$,
\[
 d_{\mathrm H}\bigl((vw)[t,t+m_i),z\bigr)
     \ge d_i m_i\ge d_*m_i.
\]
These distances count symbols in $\mathcal A$, and repeated source
words are allowed.
\end{lemma}
\begin{proof}
The finite products are bounded below by
$1-\sum_{j\ge1}\varepsilon_j=3/4$, proving that $d_*>0$.
At level one, distinct one-symbol words have distance one and the
offset assertion is vacuous.

Suppose both assertions hold at level $i$.
Distinct level-$(i+1)$ words differ in at least
$(1-\varepsilon_i)N_i$ constituent labels by
\Cref{prop:words}, Property~(\ref{word:distance}).
Each unequal constituent pair contributes at least $d_i m_i$
symbol mismatches, giving the bound $d_{i+1}m_{i+1}$.

Now consider a nonzero offset $t<m_{i+1}$ in two concatenated
level-$(i+1)$ words, compared with a third such word.
If $m_i$ divides $t$, the offset $t/m_i$ lies between $1$ and $N_i-1$.
The constituent-label offset requirement gives at least
$(1-\varepsilon_i)N_i$ unequal pairs, and the preceding distinct-word
bound again gives $d_{i+1}m_{i+1}$ symbol mismatches.
If $m_i$ does not divide $t$, each of the $N_i$ aligned target
constituents faces a length-$m_i$ source window with the same nonzero
offset modulo $m_i$.  That source window crosses exactly two
consecutive legal level-$i$ words.  The inductive offset assertion
therefore gives at least $d_i m_i$ mismatches on each target
constituent.  It applies also when the two source constituents lie
on opposite sides of the join of the two larger source words.
Adding over the disjoint target constituents gives
$d_i m_{i+1}\ge d_{i+1}m_{i+1}$, completing the induction.
\end{proof}

\subsection{The number and size of the available scales}

The word construction must retain enough label information across
levels to compete with the summable smooth-switching costs.
The following estimates are the quantitative input for that comparison.

\begin{lemma}[Word rates]\label{lem:rates}
Let $R_i=\ell_i/m_i$.  Every $m_i$ is a power of two, with
\[
  \log m_i=\sum_{j<i}K_j.
\]
The rates satisfy
\begin{equation}\label{eq:word-rate-decay}
  R_{i+1}\le\frac{\varepsilon_i}{100}R_i,
  \qquad R_i\longrightarrow0.
\end{equation}
They also have the lower bound
\begin{equation}\label{eq:word-rate-lower}
 R_i\ge R_1\,200^{-(i-1)}
          2^{-(i-1)(i+4)/2}
     \ge R_1\,2^{-(i-1)(i+20)/2},
\end{equation}
and
\begin{equation}\label{eq:word-rate-divergence}
  \sum_{i=1}^{\infty}K_iR_i^{16}=\infty.
\end{equation}
\end{lemma}
\begin{proof}
The formula for $m_i$ follows from $m_1=1$ and $N_i=2^{K_i}$.
Dividing \eqref{eq:word-rounding} by $m_{i+1}=N_i m_i$ gives
\[
 \frac{\varepsilon_i}{200}R_i
 \le R_{i+1}\le\frac{\varepsilon_i}{100}R_i.
\]
Since $\varepsilon_i\le1/8$, the upper bound implies
$R_i\le800^{-(i-1)}R_1\to0$, proving
\eqref{eq:word-rate-decay}. Iterating the lower bound and using
$\sum_{j<i}(j+2)=(i-1)(i+4)/2$ prove the first inequality in
\eqref{eq:word-rate-lower}; the second uses $\log200<8$.
Also \eqref{eq:word-growth} implies
$K_i\ge K_1^{2^{i-1}}$.  Hence
\[
 \log(K_iR_i^{16})
  \ge 2^{i-1}\log K_1+16\log R_1-8(i-1)(i+20)
  \longrightarrow\infty.
\]
In particular the terms in \eqref{eq:word-rate-divergence} do not
tend to zero, and the series diverges.
\end{proof}

The families are now fixed.  Their per-word approximation transfers
short-window entropy to any invariant measure on the resulting
sequence space, while \Cref{lem:recognition} recovers block labels
and phases from small symbol errors.  We next construct that space
and prove the required measure-theoretic statements.
